\clearpage\selectlanguage{ngerman}\fancyhead[L]{\small NACHLASS.}\pdfbookmark[0]{BESTIMMUNG UNSERER REIHE DURCH EINE DIFFERENTIALGLEICHUNG ZWEITER ORDNUNG.}{nachlass216}
\SourceLeaf{216}
\begin{center}
\textbf{NACHLASS.}\par\medskip
\textbf{BESTIMMUNG UNSERER REIHE}\par\medskip
\textbf{DURCH EINE DIFFERENTIALGLEICHUNG ZWEITER ORDNUNG.}\par\medskip
\rule{0.08\linewidth}{0.3pt}\par\bigskip
\end{center}
\SixthArticle{38}
Setzen wir der Kürze halber \(F(\alpha,\beta,\gamma,x)=P\), so haben wir \mbox{nach Artikel 4}
\[
\frac{\mathrm{d}P}{\mathrm{d}x}=\frac{\alpha\beta}{\gamma}
F(\alpha+1,\beta+1,\gamma+1,x)
\]
und durch erneutes Differenzieren daraus
\[
\frac{\mathrm{dd}P}{\mathrm{d}x^2}
=\frac{\alpha\beta(\alpha+1)(\beta+1)}{\gamma(\gamma+1)}
F(\alpha+2,\beta+2,\gamma+2,x)
\]
Damit liefert Gleichung IX aus Artikel 10
\[
\text{[80]}\qquad
0=\alpha\beta P-(\gamma-(\alpha+\beta+1)x)
\frac{\mathrm{d}P}{\mathrm{d}x}-(x-xx)\frac{\mathrm{dd}P}{\mathrm{d}x^2}
\]
Diese Differentialgleichung zweiter Ordnung kann als genauere Definition unserer Funktion angesehen werden. Da aber \(P=F(\alpha,\beta,\gamma,x)\) nur eine partikuläre und keine allgemeine Lösung ist (weil keine Konstanten hinzugekommen sind), muss man die Bedingung hinzufügen, dass \(P\) mit dem Wert 1 für \(x=0\) beginnt und zugleich für denselben Wert von \(x\)
\(\frac{\mathrm{d}P}{\mathrm{d}x}=\frac{\alpha\beta}{\gamma}\)
sowie
\(\frac{\mathrm{dd}P}{\mathrm{d}x^2}
=\frac{\alpha\beta(\alpha+1)(\beta+1)}{\gamma(\gamma+1)}\)
angesetzt werden.

Für jeden Wert von \(x\), zu dem man vom Wert \(x=0\) aus in stetigen Schritten gelangt, ohne jedoch den Wert \(x=1\) zu erreichen, für den \(x-xx=0\) ist, wird \(P\) also eine vollständig bestimmte Größe sein. Offenbar kann man aber \emph{auf diese Weise} zu positiven reellen Werten von \(x\), die größer als eins sind, nur über imaginäre Werte gelangen.\SourceLeaf{217} Da dies auf unendlich viele verschiedene Arten ohne Verletzung der Stetigkeit geschehen kann, ist noch nicht klar, ob demselben Wert von \(x\) mehrere, ja sogar unendlich viele verschiedene Werte von \(P\) entsprechen, wie es bekanntermaßen bei mehreren geläufigeren transzendenten Funktionen vorkommt. Eine ausführlichere Erörterung dieses Gegenstands behalten wir uns für später vor, da hier vor allem der Fall behandelt wird, in dem \(x\) kleiner als die positive Einheit oder zumindest nicht größer als diese genommen wird und \(P\) gleich der Summe der Reihe \(F(\alpha,\beta,\gamma,x)\) ist.

\SixthArticle{39}
Schreibt man in Gleichung 80 \(1-y\) anstelle von \(x\), so geht sie über in
\[
\begin{aligned}&0=\alpha\beta P-(\alpha+\beta+1-\gamma-(\alpha+\beta+1)y)
\frac{\mathrm{d}P}{\mathrm{d}y}\\&-(y-yy)\frac{\mathrm{dd}P}{\mathrm{d}y^2}\end{aligned}
\]
Diese hat eine ähnliche Form wie die frühere. Daraus ergibt sich sofort eine andere partikuläre Lösung:
\[
\begin{aligned}&P=F(\alpha,\beta,\alpha+\beta+1-\gamma,y)\\&=F(\alpha,\beta,\alpha+\beta+1-\gamma,1-x)\end{aligned}
\]
Nach bekannten Grundsätzen folgt daraus die allgemeine Lösung der Gleichung 80,
\[
\begin{aligned}&\text{[81]}\qquad
P=MF(\alpha,\beta,\gamma,x)\\&+NF(\alpha,\beta,\alpha+\beta+1-\gamma,1-x)\end{aligned}
\]
wobei \(M,N\) beliebige Konstanten bezeichnen.

Nebenbei bemerken wir hier noch: Auf die Form der Gleichung 80 lässt sich leicht die folgende allgemeinere Gleichung zurückführen:
\[
0=AP+(B+Cy)\frac{\mathrm{d}P}{\mathrm{d}y}
+(D+Ey+Fyy)\frac{\mathrm{dd}P}{\mathrm{d}y^2}
\]
Die Wurzeln der Gleichung \(0=D+Ey+Fyy\) seien nämlich \(y=a,y=b\), oder \(D+Ey+Fyy\) sei identisch gleich dem Produkt \(F(y-a)(y-b)\). Offenbar geht sie dann durch die Setzung \(\frac{y-a}{b-a}=x\) und die Bestimmung von \(\alpha,\beta,\gamma\) derart, dass
\[
\alpha\beta=\frac{A}{F},\quad
\alpha+\beta+1=\frac{C}{F},\quad
\gamma=-\frac{B+aC}{F(b-a)}
\]
gilt, in Gleichung 80 über.

\SixthArticle{40}
Mithilfe der Differentialgleichung 80 lassen sich mehrere sehr bemerkenswerte Sätze über unsere Reihe herleiten, teils allgemeine, teils speziellere. Wir zweifeln nicht daran, dass noch viele weitere und bedeutendere verborgen liegen, die künftigen Untersuchungen vorbehalten sind. Stellen wir hier diejenigen zusammen, die wir bisher aufdecken konnten.
\par\SourceLeaf{218}
Setzen wir $P=(1-x)^\mu P'$, so folgt
\[
\begin{aligned}
\frac{\mathrm{d}P}{\mathrm{d}x}&=-\mu(1-x)^{\mu-1}P'+(1-x)^\mu\frac{\mathrm{d}P'}{\mathrm{d}x}\\
\frac{\mathrm{dd}P}{\mathrm{d}x^2}&=\mu(\mu-1)(1-x)^{\mu-2}P'-2\mu(1-x)^{\mu-1}\frac{\mathrm{d}P'}{\mathrm{d}x}\\&\qquad+(1-x)^\mu\frac{\mathrm{dd}P'}{\mathrm{d}x^2}
\end{aligned}
\]
Setzt man diese Werte in Gleichung 80 ein und dividiert durch $(1-x)^{\mu-1}$, so erhält man
\[
\begin{aligned}
0={}&P'\{\alpha\beta(1-x)+(\gamma-(\alpha+\beta+1)x)\mu-x(\mu\mu-\mu)\}\\
&-\frac{\mathrm{d}P'}{\mathrm{d}x}\{(\gamma-(\alpha+\beta+1)x)-2\mu x\}(1-x)\\&-\frac{\mathrm{dd}P'}{\mathrm{d}x^2}\{x-xx\}(1-x)
\end{aligned}
\]
Bestimmen wir $\mu$ so, dass der Faktor von $P'$ durch $1-x$ teilbar wird. Das geschieht, wenn man entweder $\mu=0$ oder $\mu=\gamma-\alpha-\beta$ setzt. Die erste Annahme würde nichts Neues ergeben; durch Einsetzen des letzteren Wertes erhält man jedoch
\[
\begin{aligned}&0=P'\{\alpha\beta-\alpha\gamma-\beta\gamma+\gamma\gamma\}-\frac{\mathrm{d}P'}{\mathrm{d}x}\{\gamma-(2\gamma-\alpha-\beta+1)x\}\\&-\frac{\mathrm{dd}P'}{\mathrm{d}x^2}\{x-xx\}\end{aligned}
\]
oder
\[
\begin{aligned}&0=P'(\gamma-\alpha)(\gamma-\beta)-\frac{\mathrm{d}P'}{\mathrm{d}x}\{\gamma-((\gamma-\alpha)+(\gamma-\beta)+1)x\}\\&-\frac{\mathrm{dd}P'}{\mathrm{d}x^2}(x-xx)\end{aligned}
\]
Dies hat genau dieselbe Form wie Gleichung 80. Da nun für $x=0$ offenbar $P'=1$ und $\frac{\mathrm{d}P'}{\mathrm{d}x}=\frac{\alpha\beta}{\gamma}-\mu=\frac{(\gamma-\alpha)(\gamma-\beta)}{\gamma}$ gelten, ist die Lösung dieser Gleichung ersichtlich $P'=F(\gamma-\alpha,\gamma-\beta,\gamma,x)$. Somit gilt allgemein
\[
\text{[82]}\qquad F(\gamma-\alpha,\gamma-\beta,\gamma,x)=(1-x)^{\alpha+\beta-\gamma}F(\alpha,\beta,\gamma,x)
\]
Hieraus ist die Umformung der Reihe
\[
1+\frac{2\,.\,8}{9}x+\frac{3\,.\,8\,.\,10}{9\,.\,11}xx+\frac{4\,.\,8\,.\,10\,.\,12}{9\,.\,11\,.\,13}x^3+\text{usw.}=F(2,4,\tfrac{9}{2},x)
\]
in
\[
\begin{aligned}&(1-x)^{-\frac{3}{2}}\left(1+\frac{1\,.\,5}{2\,.\,9}x+\frac{1\,.\,3\,.\,5\,.\,7}{2\,.\,4\,.\,9\,.\,11}xx+\text{usw.}\right)\\&=(1-x)^{-\frac{3}{2}}F(\tfrac{5}{2},\tfrac{1}{2},\tfrac{9}{2},x)\end{aligned}
\]
herzuleiten, die wir in den Berliner Astronomischen Ephemeriden für 1814, S. 257 [Zusatz zu Art. 90 und 100 der Theoria motus] ohne Beweis angegeben hatten.

